% Generated draft: requires mathematical and source review.
\documentclass[11pt]{article}
\usepackage[margin=25mm]{geometry}
\usepackage{fontspec}
\setmainfont{Latin Modern Roman}
\usepackage{amsmath,amssymb,mathtools,unicode-math}
\setmathfont{Latin Modern Math}
\usepackage{enumitem,needspace,xurl,hyperref}
\setlength{\parindent}{0pt}
\setlength{\parskip}{5pt}
\setlength{\emergencystretch}{3em}
\newcommand{\N}{\mathbb{N}}
\newcommand{\Z}{\mathbb{Z}}
\newcommand{\Q}{\mathbb{Q}}
\newcommand{\R}{\mathbb{R}}
\newcommand{\C}{\mathbb{C}}
\newcommand{\F}{\mathbb{F}}
\newcommand{\A}{\mathbb{A}}
\newcommand{\PP}{\mathbb{P}}
\newcommand{\E}{\mathbb{E}}
\newcommand{\eps}{\varepsilon}
\newcommand{\dd}{\,\mathrm{d}}
\begin{document}
\section*{UnsolvedMath ID: 30003700}
\textbf{Source problem code:} catalogue entry
\subsection{Definitions and mathematical statement}
Write a monic polynomial of degree \(n\) as
\[
P(z)=z^n+a_1z^{n-1}+\cdots+a_n.
\]
The symbols \(a_1,\ldots,a_n\) denote its coefficients. A root is an element \(r\) of an appropriate ambient domain satisfying \(P(r)=0\).

\textbf{Source problem.} Find and understand formulas for the roots of \(P\) in terms of its coefficients \(a_1,\ldots,a_n\).

\textbf{Unresolved qualifications.} The source does not specify the coefficient domain, the ambient domain containing the roots, the intended range of degrees, or the operations and functions permitted in a formula. It also does not specify whether formulas must describe every root with multiplicity, how branches are to be chosen, whether repeated roots are included, or whether formulas may be restricted to generic coefficients. The phrase ``understand formulas'' supplies no precise criterion for success. Thus the supplied request is a research objective rather than a fully quantified conjecture; a more specific mathematical assertion cannot be extracted without additional choices.
\subsection{Short English statement}
Seek formulas expressing polynomial roots in terms of their coefficients, together with an explanation of those formulas. The source leaves the domains, permitted functions, and scope of the requested formulas unspecified. Its supplied status is \texttt{partially\_solved}; this draft does not independently verify that assessment.
\subsection{Sources}
\begin{itemize}
\item Local UnsolvedMath registry: supplied complete source record, permanent ID \(30003700\), title \emph{Coefficient Formulas for Polynomial Roots}, exact source problem code \texttt{catalogue entry}, collection \emph{Oberwolfach Reports Open Problems}. The proposed year is recorded as \(2018\); no proposer is named.
\item Supplied source citation: \emph{Topology of Arrangements and Representation Stability (2019)}, Oberwolfach Reports, pp.~43--123, DOI \texttt{10.4171/owr/2018/2}. \url{https://doi.org/10.4171/owr/2018/2}. The background identifies Oberwolfach Reports as reports of mathematical workshops and their open-problem discussions. The citation is reproduced as supplied; its bibliographic details have not been independently verified.
\item Supplied literature source: Farb and Wolfson, resolvent degree. \url{https://arxiv.org/abs/1803.04063}.
\item Literature assessment dated exactly \texttt{2026-08-21}: supplied dated curation, not independently verified. The record labels the problem partially solved and states that the impossibility of radical formulas from degree five onward and classical reductions such as Bring's are established, while the broader request for formulas and understanding of roots remains incompletely resolved. The background describes this as dated literature triage derived from the supplied OWR curation and explicitly requests independent verification of time-sensitive status and cited claims.
\end{itemize}
\end{document}
